% ============================================================
%  bundle1.tex — Abstract + Introduction + Background
%  % ============================================================
%  bundle1.tex — Abstract + Introduction + Background
%  % ============================================================
%  bundle1.tex — Abstract + Introduction + Background
%  % ============================================================
%  bundle1.tex — Abstract + Introduction + Background
%  \input{sections/bundle1} from main.tex
% ============================================================

\begin{abstract}
India's Defence and critical-infrastructure surveillance network spans thousands of kilometers, with critical information scattered across unconnected sensors and surveillance systems. This fragmentation, along with human operation for threat detection and classification, makes the system suffer from high false-alarm rates, which constrain timely and reliable threat assessment.
Recent advancements in technology, AI, and Machine Learning offer new methods to enhance the system, such as multi-sensor fusion and AI analysis for automating surveillance systems, leaving only the critical decision-making part with manual operators. This survey provides a structured review of the five major domains of such an AI-based system that can potentially automate surveillance and reduce operators' fatigue and high false-alarm rates, providing accurate threat detection and classification. Our principal observation is that despite the growth in AI and ML, there exists no such framework that unifies diverse sensors and performs accurate threat detection and five-level threat classification across multiple domains. These gaps motivate the design and study of a unified, multi-agent, AI-driven surveillance platform tailored to Indian Defence requirements.
\end{abstract}

\vspace{\baselineskip}

\begin{IEEEkeywords}
Defence surveillance, multi-modal sensor fusion, multi-agent systems, threat detection and classification, explainable artificial intelligence, human-in-the-loop, adversarial robustness, edge computing, deep learning, literature survey
\end{IEEEkeywords}

\section{Introduction}
\label{sec:introduction}

India maintains active surveillance obligations along more than 15{,}000~km of land borders with Pakistan, Bangladesh, Myanmar, and Nepal, a 7{,}500~km coastline, and an expanding portfolio of critical infrastructure installations~\cite{indiaai2024springer}. Each segment employs its own constellation of radars, thermal imagers, unattended ground sensors, electro-optical cameras, and communication relays, generating data volumes that far exceed the analytical capacity of human operators~\cite{meng2023multisource}. Although individual sensor technologies have matured significantly over the years, the lack of a unified system that can connect and reason across multiple sensor types and geographic regions leaves today's surveillance infrastructure fundamentally fragile. A scientometric analysis of defence-technology publications between 2015 and 2025 confirms a fourfold increase in AI-related land-combat research, yet integrated multi-domain surveillance architectures remain conspicuously under-represented~\cite{emerging2025scientometric}.

Two critical problems make this gap even more dangerous. First, the high false alarm rates seen across deployed perimeter security systems place a growing burden on human operators, every false trigger pulls attention and resources away from threats that actually matter~\cite{deepsurvey2023anomaly}. Second, most existing alarm systems treat all alerts the same way, meaning a deer crossing a perimeter sensor and a coordinated military incursion receive the same level of urgency, making it impossible for operators to prioritise their response effectively~\cite{cctv2024anomaly}. Decision-support systems that incorporate explainability techniques such as SHAP values or attention-based saliency maps have shown encouraging results in civilian surveillance for identifying and triaging anomalies~\cite{zhang2022xai}, but translating these approaches into military environments, where decisions must be made faster, accountability is higher, and the consequences of errors are far greater, remains an unsolved challenge~\cite{military2024xai}.

Multi-agent AI architectures offer a principled response to both failures. Rather than replacing human judgment, a network of specialised software agents can function as a filtering and escalation layer that operationalises algorithmic situational awareness: each agent ingests a defined sensor stream, applies domain-specific inference, and feeds a threat-correlation agent that synthesises inputs into a graduated alert~\cite{hassan2025multiagent, agenticcyber2025}. Such architectures have demonstrated measurable gains in cyber-threat correlation and battlefield command support~\cite{zhao2022deepai}, and purpose-built human-autonomy teaming frameworks preserve commander authority over critical engagement decisions~\cite{autonomyteaming2024, meaningful2024hitl}. Extending this paradigm to physical-domain defence surveillance requires integration across five technical pillars that, to date, have been investigated largely in isolation.

This paper makes two contributions. First, it provides a structured literature review spanning multi-modal sensor fusion~\cite{meng2023multisource}, multi-agent threat detection~\cite{hassan2025multiagent, agenticcyber2025}, false-alarm reduction and explainable AI~\cite{zhang2022xai, deepsurvey2023anomaly}, geospatial awareness with human-in-the-loop control~\cite{karimipour2021iot, hitl2023trust}, and real-time distributed architectures with adversarial robustness~\cite{flink2024kafka, advyolo2022}. Second, it identifies and formalises the integration gap across these five domains as the primary bottleneck obstructing India's transition from sensor-rich but intelligence-poor surveillance to a unified, graduated threat-management grid. Section~II traces the historical and policy background. Sections~III through~VII examine each domain. Section~VIII presents a comparative taxonomy, analyses open gaps, outlines a proposed architectural direction, and concludes the paper.



\section{Background}
\label{sec:background}

India's post-independence approach towards managing its borders can be characterized by manpower-centric operations relying heavily on traditional techniques such as fences, patrolling, optical lookout posts, and high-frequency radio communications~\cite{indiaai2024springer}. Over the last several decades, the primary method for collecting information about the conditions at the frontline was through the eyes and ears of the soldiers stationed in border posts combined with basic surveillance tools such as binoculars, alarm wires, and searchlight systems installed on vehicles~\cite{indiaai2024springer}. Though a known approach, this technique does not scale effectively to India's size, variety of geographical regions, and harsh weather conditions and provides very little ability to conduct ongoing surveillance around-the-clock.

The Comprehensive Integrated Border Management System (CIBMS) may be considered as the largest breakthrough in terms of both technology and doctrine of Indian border surveillance. Designed to replace traditional physical barriers along the India-Pakistan and India-Bangladesh borders, CIBMS consists of various sensors ranging from thermal imaging cameras and radar systems mounted on tethered aerostats, to laser and infrared sensors that detect intrusion attempts and unattended ground sensors. Further refinements added various radars for detecting battlefield threats, human presence, and seismic detectors able to pick up sounds created during tunnel construction~\cite{indiaai2024springer}. Despite the transition to C2-centric surveillance, CIBMS' alert generation still relies on rule-based threshold logic and does not discriminate between false positives and actual violations of border integrity~\cite{cctv2024anomaly}.

There is a parallel revolution underway in the military surveillance R\&D sphere, with rule-based analysis giving way to machine learning and deep neural networks allowing for processing vast amounts of sensor data faster than even human analysts are able to~\cite{threatdetect2024rt}. Several convolutional architectures were shown to have good performance at detecting military vehicles and weapon systems~\cite{weapon2023yolov8} as well as military equipment visible in satellite imagery~\cite{satellite2024yolov8}. Combined CNN-LSTM architectures for processing sequential data improved anomaly detection accuracy and decreased false-positive rates~\cite{deepsurvey2023anomaly}. An example of successful implementation of AI pipelines is Project Maven launched by the U.S.\ Department of Defense, it showed that with machine learning assistance, 20 people were able to process the same amount of full-motion video footage as 2{,}000 people could without ML~\cite{cset2024maven, rand2023maven}.

The response from India's defence leadership has been quite proactive. The Ministry of Defence prioritized development of artificial intelligence solutions and increased CAIR's remit to include autonomous surveillance systems, target recognition, and AI-powered command systems~\cite{indiaai2024springer}. Commercial off-the-shelf AI providers and data fusion systems have previously been assessed for cross-database correlations and anomaly detection capability in intelligence analysis~\cite{cset2025ai}. Nevertheless, an analysis of India's defence-AI ecosystem reveals a prevalence of pilots, which are stovepipe within each branch of the defence forces due to insufficient interoperability capabilities and conceptually divergent data semantics~\cite{indiaai2024springer}.

There are several key concerns that drive this survey effort. One of the biggest problems is bandwidth limitations in contested border regions, which requires development of edge-computing solutions able to provide preliminary threat assessment at sensor nodes~\cite{flink2024kafka}. In addition, environmental noise sources, including sandstorms occurring near the border with Rajasthan and river level changes caused by seasonal monsoons in the border region with Bangladesh, generate long-term false alarm signals, which overwhelm threshold-based detection methods~\cite{deepsurvey2023anomaly}. Furthermore, deep learning classifiers suffer from poor interpretability, raising concerns about military accountability when making decisions related to escalation of force in accordance with international humanitarian laws~\cite{military2024xai, meaningful2024hitl}. Legacy sensor devices used in the CIBMS sector lack unified data interfaces, making real-time fusion with AI-enabled sensors difficult~\cite{meng2023multisource}. Finally, modern adversaries are able to attack object detection models via adversarial patches and GPS signal spoofing~\cite{advyolo2022, milai2025risk}. All of the mentioned concerns represent open challenges that this survey seeks to address, starting with a systematic overview of the five key pillars.

% NOTE: table* used for full-width layout; 5 content columns
%       cannot fit legibly within a single IEEEtran column.
\begin{table*}[!t]
\centering
\caption{Taxonomy of Defence Surveillance Literature}
\label{tab:taxonomy}
\renewcommand{\arraystretch}{1.3}
\footnotesize
\begin{tabular}{>{\raggedright\arraybackslash}p{1.5cm}>{\raggedright\arraybackslash}p{3.6cm}>{\raggedright\arraybackslash}p{2.4cm}>{\raggedright\arraybackslash}p{2.6cm}>{\raggedright\arraybackslash}p{3.8cm}}
\hline
\textbf{Domain} & \textbf{Representative Approaches} & \textbf{Sensor / Data Modality} & \textbf{Key Technique} & \textbf{Limitation Noted} \\
\hline
A: Sensor Fusion
  & Multi-source information fusion for C3I systems; attention-based cross-modal alignment
  & Radar, thermal, acoustic, optical, RF
  & Feature-level and decision-level multi-source fusion
  & Evaluated on pairwise modalities only; no adversarial degradation testing \\
\hline
B: Multi-Agent AI
  & Modular MAS with role-specialised sub-agents; YOLOv8 military object detection
  & Video, radar returns, satellite imagery, system logs
  & Agent decomposition with cross-source threat correlation
  & Homogeneous task distribution; no cross-modal agent orchestration \\
\hline
C: False Alarm / XAI
  & Deep-learning anomaly detection; SHAP/LIME explanation methods
  & Surveillance video, CCTV feeds, IDS telemetry
  & Temporal consistency filtering; post-hoc feature attribution
  & XAI latency not evaluated within real-time alert pipelines \\
\hline
D: Geospatial / HITL
  & GIS-based threat dashboards; trust calibration for autonomy teaming
  & GIS terrain layers, sensor metadata, force positions
  & Terrain-aware threat mapping; human authorisation loops
  & Tested in single-operator, single-threat scenarios only \\
\hline
E: Distributed / Adversarial
  & Kafka + Flink real-time streaming; adversarial patch attacks on YOLO detectors
  & IoT telemetry, edge video, RF uplinks
  & Edge preprocessing with stream-processing inference; adversarial training
  & No evaluation under simultaneous node failure and adversarial injection \\
\hline
\end{tabular}
\end{table*}
 from main.tex
% ============================================================

\begin{abstract}
India's Defence and critical-infrastructure surveillance network spans thousands of kilometers, with critical information scattered across unconnected sensors and surveillance systems. This fragmentation, along with human operation for threat detection and classification, makes the system suffer from high false-alarm rates, which constrain timely and reliable threat assessment.
Recent advancements in technology, AI, and Machine Learning offer new methods to enhance the system, such as multi-sensor fusion and AI analysis for automating surveillance systems, leaving only the critical decision-making part with manual operators. This survey provides a structured review of the five major domains of such an AI-based system that can potentially automate surveillance and reduce operators' fatigue and high false-alarm rates, providing accurate threat detection and classification. Our principal observation is that despite the growth in AI and ML, there exists no such framework that unifies diverse sensors and performs accurate threat detection and five-level threat classification across multiple domains. These gaps motivate the design and study of a unified, multi-agent, AI-driven surveillance platform tailored to Indian Defence requirements.
\end{abstract}

\vspace{\baselineskip}

\begin{IEEEkeywords}
Defence surveillance, multi-modal sensor fusion, multi-agent systems, threat detection and classification, explainable artificial intelligence, human-in-the-loop, adversarial robustness, edge computing, deep learning, literature survey
\end{IEEEkeywords}

\section{Introduction}
\label{sec:introduction}

India maintains active surveillance obligations along more than 15{,}000~km of land borders with Pakistan, Bangladesh, Myanmar, and Nepal, a 7{,}500~km coastline, and an expanding portfolio of critical infrastructure installations~\cite{indiaai2024springer}. Each segment employs its own constellation of radars, thermal imagers, unattended ground sensors, electro-optical cameras, and communication relays, generating data volumes that far exceed the analytical capacity of human operators~\cite{meng2023multisource}. Although individual sensor technologies have matured significantly over the years, the lack of a unified system that can connect and reason across multiple sensor types and geographic regions leaves today's surveillance infrastructure fundamentally fragile. A scientometric analysis of defence-technology publications between 2015 and 2025 confirms a fourfold increase in AI-related land-combat research, yet integrated multi-domain surveillance architectures remain conspicuously under-represented~\cite{emerging2025scientometric}.

Two critical problems make this gap even more dangerous. First, the high false alarm rates seen across deployed perimeter security systems place a growing burden on human operators, every false trigger pulls attention and resources away from threats that actually matter~\cite{deepsurvey2023anomaly}. Second, most existing alarm systems treat all alerts the same way, meaning a deer crossing a perimeter sensor and a coordinated military incursion receive the same level of urgency, making it impossible for operators to prioritise their response effectively~\cite{cctv2024anomaly}. Decision-support systems that incorporate explainability techniques such as SHAP values or attention-based saliency maps have shown encouraging results in civilian surveillance for identifying and triaging anomalies~\cite{zhang2022xai}, but translating these approaches into military environments, where decisions must be made faster, accountability is higher, and the consequences of errors are far greater, remains an unsolved challenge~\cite{military2024xai}.

Multi-agent AI architectures offer a principled response to both failures. Rather than replacing human judgment, a network of specialised software agents can function as a filtering and escalation layer that operationalises algorithmic situational awareness: each agent ingests a defined sensor stream, applies domain-specific inference, and feeds a threat-correlation agent that synthesises inputs into a graduated alert~\cite{hassan2025multiagent, agenticcyber2025}. Such architectures have demonstrated measurable gains in cyber-threat correlation and battlefield command support~\cite{zhao2022deepai}, and purpose-built human-autonomy teaming frameworks preserve commander authority over critical engagement decisions~\cite{autonomyteaming2024, meaningful2024hitl}. Extending this paradigm to physical-domain defence surveillance requires integration across five technical pillars that, to date, have been investigated largely in isolation.

This paper makes two contributions. First, it provides a structured literature review spanning multi-modal sensor fusion~\cite{meng2023multisource}, multi-agent threat detection~\cite{hassan2025multiagent, agenticcyber2025}, false-alarm reduction and explainable AI~\cite{zhang2022xai, deepsurvey2023anomaly}, geospatial awareness with human-in-the-loop control~\cite{karimipour2021iot, hitl2023trust}, and real-time distributed architectures with adversarial robustness~\cite{flink2024kafka, advyolo2022}. Second, it identifies and formalises the integration gap across these five domains as the primary bottleneck obstructing India's transition from sensor-rich but intelligence-poor surveillance to a unified, graduated threat-management grid. Section~II traces the historical and policy background. Sections~III through~VII examine each domain. Section~VIII presents a comparative taxonomy, analyses open gaps, outlines a proposed architectural direction, and concludes the paper.



\section{Background}
\label{sec:background}

India's post-independence approach towards managing its borders can be characterized by manpower-centric operations relying heavily on traditional techniques such as fences, patrolling, optical lookout posts, and high-frequency radio communications~\cite{indiaai2024springer}. Over the last several decades, the primary method for collecting information about the conditions at the frontline was through the eyes and ears of the soldiers stationed in border posts combined with basic surveillance tools such as binoculars, alarm wires, and searchlight systems installed on vehicles~\cite{indiaai2024springer}. Though a known approach, this technique does not scale effectively to India's size, variety of geographical regions, and harsh weather conditions and provides very little ability to conduct ongoing surveillance around-the-clock.

The Comprehensive Integrated Border Management System (CIBMS) may be considered as the largest breakthrough in terms of both technology and doctrine of Indian border surveillance. Designed to replace traditional physical barriers along the India-Pakistan and India-Bangladesh borders, CIBMS consists of various sensors ranging from thermal imaging cameras and radar systems mounted on tethered aerostats, to laser and infrared sensors that detect intrusion attempts and unattended ground sensors. Further refinements added various radars for detecting battlefield threats, human presence, and seismic detectors able to pick up sounds created during tunnel construction~\cite{indiaai2024springer}. Despite the transition to C2-centric surveillance, CIBMS' alert generation still relies on rule-based threshold logic and does not discriminate between false positives and actual violations of border integrity~\cite{cctv2024anomaly}.

There is a parallel revolution underway in the military surveillance R\&D sphere, with rule-based analysis giving way to machine learning and deep neural networks allowing for processing vast amounts of sensor data faster than even human analysts are able to~\cite{threatdetect2024rt}. Several convolutional architectures were shown to have good performance at detecting military vehicles and weapon systems~\cite{weapon2023yolov8} as well as military equipment visible in satellite imagery~\cite{satellite2024yolov8}. Combined CNN-LSTM architectures for processing sequential data improved anomaly detection accuracy and decreased false-positive rates~\cite{deepsurvey2023anomaly}. An example of successful implementation of AI pipelines is Project Maven launched by the U.S.\ Department of Defense, it showed that with machine learning assistance, 20 people were able to process the same amount of full-motion video footage as 2{,}000 people could without ML~\cite{cset2024maven, rand2023maven}.

The response from India's defence leadership has been quite proactive. The Ministry of Defence prioritized development of artificial intelligence solutions and increased CAIR's remit to include autonomous surveillance systems, target recognition, and AI-powered command systems~\cite{indiaai2024springer}. Commercial off-the-shelf AI providers and data fusion systems have previously been assessed for cross-database correlations and anomaly detection capability in intelligence analysis~\cite{cset2025ai}. Nevertheless, an analysis of India's defence-AI ecosystem reveals a prevalence of pilots, which are stovepipe within each branch of the defence forces due to insufficient interoperability capabilities and conceptually divergent data semantics~\cite{indiaai2024springer}.

There are several key concerns that drive this survey effort. One of the biggest problems is bandwidth limitations in contested border regions, which requires development of edge-computing solutions able to provide preliminary threat assessment at sensor nodes~\cite{flink2024kafka}. In addition, environmental noise sources, including sandstorms occurring near the border with Rajasthan and river level changes caused by seasonal monsoons in the border region with Bangladesh, generate long-term false alarm signals, which overwhelm threshold-based detection methods~\cite{deepsurvey2023anomaly}. Furthermore, deep learning classifiers suffer from poor interpretability, raising concerns about military accountability when making decisions related to escalation of force in accordance with international humanitarian laws~\cite{military2024xai, meaningful2024hitl}. Legacy sensor devices used in the CIBMS sector lack unified data interfaces, making real-time fusion with AI-enabled sensors difficult~\cite{meng2023multisource}. Finally, modern adversaries are able to attack object detection models via adversarial patches and GPS signal spoofing~\cite{advyolo2022, milai2025risk}. All of the mentioned concerns represent open challenges that this survey seeks to address, starting with a systematic overview of the five key pillars.

% NOTE: table* used for full-width layout; 5 content columns
%       cannot fit legibly within a single IEEEtran column.
\begin{table*}[!t]
\centering
\caption{Taxonomy of Defence Surveillance Literature}
\label{tab:taxonomy}
\renewcommand{\arraystretch}{1.3}
\footnotesize
\begin{tabular}{>{\raggedright\arraybackslash}p{1.5cm}>{\raggedright\arraybackslash}p{3.6cm}>{\raggedright\arraybackslash}p{2.4cm}>{\raggedright\arraybackslash}p{2.6cm}>{\raggedright\arraybackslash}p{3.8cm}}
\hline
\textbf{Domain} & \textbf{Representative Approaches} & \textbf{Sensor / Data Modality} & \textbf{Key Technique} & \textbf{Limitation Noted} \\
\hline
A: Sensor Fusion
  & Multi-source information fusion for C3I systems; attention-based cross-modal alignment
  & Radar, thermal, acoustic, optical, RF
  & Feature-level and decision-level multi-source fusion
  & Evaluated on pairwise modalities only; no adversarial degradation testing \\
\hline
B: Multi-Agent AI
  & Modular MAS with role-specialised sub-agents; YOLOv8 military object detection
  & Video, radar returns, satellite imagery, system logs
  & Agent decomposition with cross-source threat correlation
  & Homogeneous task distribution; no cross-modal agent orchestration \\
\hline
C: False Alarm / XAI
  & Deep-learning anomaly detection; SHAP/LIME explanation methods
  & Surveillance video, CCTV feeds, IDS telemetry
  & Temporal consistency filtering; post-hoc feature attribution
  & XAI latency not evaluated within real-time alert pipelines \\
\hline
D: Geospatial / HITL
  & GIS-based threat dashboards; trust calibration for autonomy teaming
  & GIS terrain layers, sensor metadata, force positions
  & Terrain-aware threat mapping; human authorisation loops
  & Tested in single-operator, single-threat scenarios only \\
\hline
E: Distributed / Adversarial
  & Kafka + Flink real-time streaming; adversarial patch attacks on YOLO detectors
  & IoT telemetry, edge video, RF uplinks
  & Edge preprocessing with stream-processing inference; adversarial training
  & No evaluation under simultaneous node failure and adversarial injection \\
\hline
\end{tabular}
\end{table*}
 from main.tex
% ============================================================

\begin{abstract}
India's Defence and critical-infrastructure surveillance network spans thousands of kilometers, with critical information scattered across unconnected sensors and surveillance systems. This fragmentation, along with human operation for threat detection and classification, makes the system suffer from high false-alarm rates, which constrain timely and reliable threat assessment.
Recent advancements in technology, AI, and Machine Learning offer new methods to enhance the system, such as multi-sensor fusion and AI analysis for automating surveillance systems, leaving only the critical decision-making part with manual operators. This survey provides a structured review of the five major domains of such an AI-based system that can potentially automate surveillance and reduce operators' fatigue and high false-alarm rates, providing accurate threat detection and classification. Our principal observation is that despite the growth in AI and ML, there exists no such framework that unifies diverse sensors and performs accurate threat detection and five-level threat classification across multiple domains. These gaps motivate the design and study of a unified, multi-agent, AI-driven surveillance platform tailored to Indian Defence requirements.
\end{abstract}

\vspace{\baselineskip}

\begin{IEEEkeywords}
Defence surveillance, multi-modal sensor fusion, multi-agent systems, threat detection and classification, explainable artificial intelligence, human-in-the-loop, adversarial robustness, edge computing, deep learning, literature survey
\end{IEEEkeywords}

\section{Introduction}
\label{sec:introduction}

India maintains active surveillance obligations along more than 15{,}000~km of land borders with Pakistan, Bangladesh, Myanmar, and Nepal, a 7{,}500~km coastline, and an expanding portfolio of critical infrastructure installations~\cite{indiaai2024springer}. Each segment employs its own constellation of radars, thermal imagers, unattended ground sensors, electro-optical cameras, and communication relays, generating data volumes that far exceed the analytical capacity of human operators~\cite{meng2023multisource}. Although individual sensor technologies have matured significantly over the years, the lack of a unified system that can connect and reason across multiple sensor types and geographic regions leaves today's surveillance infrastructure fundamentally fragile. A scientometric analysis of defence-technology publications between 2015 and 2025 confirms a fourfold increase in AI-related land-combat research, yet integrated multi-domain surveillance architectures remain conspicuously under-represented~\cite{emerging2025scientometric}.

Two critical problems make this gap even more dangerous. First, the high false alarm rates seen across deployed perimeter security systems place a growing burden on human operators, every false trigger pulls attention and resources away from threats that actually matter~\cite{deepsurvey2023anomaly}. Second, most existing alarm systems treat all alerts the same way, meaning a deer crossing a perimeter sensor and a coordinated military incursion receive the same level of urgency, making it impossible for operators to prioritise their response effectively~\cite{cctv2024anomaly}. Decision-support systems that incorporate explainability techniques such as SHAP values or attention-based saliency maps have shown encouraging results in civilian surveillance for identifying and triaging anomalies~\cite{zhang2022xai}, but translating these approaches into military environments, where decisions must be made faster, accountability is higher, and the consequences of errors are far greater, remains an unsolved challenge~\cite{military2024xai}.

Multi-agent AI architectures offer a principled response to both failures. Rather than replacing human judgment, a network of specialised software agents can function as a filtering and escalation layer that operationalises algorithmic situational awareness: each agent ingests a defined sensor stream, applies domain-specific inference, and feeds a threat-correlation agent that synthesises inputs into a graduated alert~\cite{hassan2025multiagent, agenticcyber2025}. Such architectures have demonstrated measurable gains in cyber-threat correlation and battlefield command support~\cite{zhao2022deepai}, and purpose-built human-autonomy teaming frameworks preserve commander authority over critical engagement decisions~\cite{autonomyteaming2024, meaningful2024hitl}. Extending this paradigm to physical-domain defence surveillance requires integration across five technical pillars that, to date, have been investigated largely in isolation.

This paper makes two contributions. First, it provides a structured literature review spanning multi-modal sensor fusion~\cite{meng2023multisource}, multi-agent threat detection~\cite{hassan2025multiagent, agenticcyber2025}, false-alarm reduction and explainable AI~\cite{zhang2022xai, deepsurvey2023anomaly}, geospatial awareness with human-in-the-loop control~\cite{karimipour2021iot, hitl2023trust}, and real-time distributed architectures with adversarial robustness~\cite{flink2024kafka, advyolo2022}. Second, it identifies and formalises the integration gap across these five domains as the primary bottleneck obstructing India's transition from sensor-rich but intelligence-poor surveillance to a unified, graduated threat-management grid. Section~II traces the historical and policy background. Sections~III through~VII examine each domain. Section~VIII presents a comparative taxonomy, analyses open gaps, outlines a proposed architectural direction, and concludes the paper.



\section{Background}
\label{sec:background}

India's post-independence approach towards managing its borders can be characterized by manpower-centric operations relying heavily on traditional techniques such as fences, patrolling, optical lookout posts, and high-frequency radio communications~\cite{indiaai2024springer}. Over the last several decades, the primary method for collecting information about the conditions at the frontline was through the eyes and ears of the soldiers stationed in border posts combined with basic surveillance tools such as binoculars, alarm wires, and searchlight systems installed on vehicles~\cite{indiaai2024springer}. Though a known approach, this technique does not scale effectively to India's size, variety of geographical regions, and harsh weather conditions and provides very little ability to conduct ongoing surveillance around-the-clock.

The Comprehensive Integrated Border Management System (CIBMS) may be considered as the largest breakthrough in terms of both technology and doctrine of Indian border surveillance. Designed to replace traditional physical barriers along the India-Pakistan and India-Bangladesh borders, CIBMS consists of various sensors ranging from thermal imaging cameras and radar systems mounted on tethered aerostats, to laser and infrared sensors that detect intrusion attempts and unattended ground sensors. Further refinements added various radars for detecting battlefield threats, human presence, and seismic detectors able to pick up sounds created during tunnel construction~\cite{indiaai2024springer}. Despite the transition to C2-centric surveillance, CIBMS' alert generation still relies on rule-based threshold logic and does not discriminate between false positives and actual violations of border integrity~\cite{cctv2024anomaly}.

There is a parallel revolution underway in the military surveillance R\&D sphere, with rule-based analysis giving way to machine learning and deep neural networks allowing for processing vast amounts of sensor data faster than even human analysts are able to~\cite{threatdetect2024rt}. Several convolutional architectures were shown to have good performance at detecting military vehicles and weapon systems~\cite{weapon2023yolov8} as well as military equipment visible in satellite imagery~\cite{satellite2024yolov8}. Combined CNN-LSTM architectures for processing sequential data improved anomaly detection accuracy and decreased false-positive rates~\cite{deepsurvey2023anomaly}. An example of successful implementation of AI pipelines is Project Maven launched by the U.S.\ Department of Defense, it showed that with machine learning assistance, 20 people were able to process the same amount of full-motion video footage as 2{,}000 people could without ML~\cite{cset2024maven, rand2023maven}.

The response from India's defence leadership has been quite proactive. The Ministry of Defence prioritized development of artificial intelligence solutions and increased CAIR's remit to include autonomous surveillance systems, target recognition, and AI-powered command systems~\cite{indiaai2024springer}. Commercial off-the-shelf AI providers and data fusion systems have previously been assessed for cross-database correlations and anomaly detection capability in intelligence analysis~\cite{cset2025ai}. Nevertheless, an analysis of India's defence-AI ecosystem reveals a prevalence of pilots, which are stovepipe within each branch of the defence forces due to insufficient interoperability capabilities and conceptually divergent data semantics~\cite{indiaai2024springer}.

There are several key concerns that drive this survey effort. One of the biggest problems is bandwidth limitations in contested border regions, which requires development of edge-computing solutions able to provide preliminary threat assessment at sensor nodes~\cite{flink2024kafka}. In addition, environmental noise sources, including sandstorms occurring near the border with Rajasthan and river level changes caused by seasonal monsoons in the border region with Bangladesh, generate long-term false alarm signals, which overwhelm threshold-based detection methods~\cite{deepsurvey2023anomaly}. Furthermore, deep learning classifiers suffer from poor interpretability, raising concerns about military accountability when making decisions related to escalation of force in accordance with international humanitarian laws~\cite{military2024xai, meaningful2024hitl}. Legacy sensor devices used in the CIBMS sector lack unified data interfaces, making real-time fusion with AI-enabled sensors difficult~\cite{meng2023multisource}. Finally, modern adversaries are able to attack object detection models via adversarial patches and GPS signal spoofing~\cite{advyolo2022, milai2025risk}. All of the mentioned concerns represent open challenges that this survey seeks to address, starting with a systematic overview of the five key pillars.

% NOTE: table* used for full-width layout; 5 content columns
%       cannot fit legibly within a single IEEEtran column.
\begin{table*}[!t]
\centering
\caption{Taxonomy of Defence Surveillance Literature}
\label{tab:taxonomy}
\renewcommand{\arraystretch}{1.3}
\footnotesize
\begin{tabular}{>{\raggedright\arraybackslash}p{1.5cm}>{\raggedright\arraybackslash}p{3.6cm}>{\raggedright\arraybackslash}p{2.4cm}>{\raggedright\arraybackslash}p{2.6cm}>{\raggedright\arraybackslash}p{3.8cm}}
\hline
\textbf{Domain} & \textbf{Representative Approaches} & \textbf{Sensor / Data Modality} & \textbf{Key Technique} & \textbf{Limitation Noted} \\
\hline
A: Sensor Fusion
  & Multi-source information fusion for C3I systems; attention-based cross-modal alignment
  & Radar, thermal, acoustic, optical, RF
  & Feature-level and decision-level multi-source fusion
  & Evaluated on pairwise modalities only; no adversarial degradation testing \\
\hline
B: Multi-Agent AI
  & Modular MAS with role-specialised sub-agents; YOLOv8 military object detection
  & Video, radar returns, satellite imagery, system logs
  & Agent decomposition with cross-source threat correlation
  & Homogeneous task distribution; no cross-modal agent orchestration \\
\hline
C: False Alarm / XAI
  & Deep-learning anomaly detection; SHAP/LIME explanation methods
  & Surveillance video, CCTV feeds, IDS telemetry
  & Temporal consistency filtering; post-hoc feature attribution
  & XAI latency not evaluated within real-time alert pipelines \\
\hline
D: Geospatial / HITL
  & GIS-based threat dashboards; trust calibration for autonomy teaming
  & GIS terrain layers, sensor metadata, force positions
  & Terrain-aware threat mapping; human authorisation loops
  & Tested in single-operator, single-threat scenarios only \\
\hline
E: Distributed / Adversarial
  & Kafka + Flink real-time streaming; adversarial patch attacks on YOLO detectors
  & IoT telemetry, edge video, RF uplinks
  & Edge preprocessing with stream-processing inference; adversarial training
  & No evaluation under simultaneous node failure and adversarial injection \\
\hline
\end{tabular}
\end{table*}
 from main.tex
% ============================================================

\begin{abstract}
India's Defence and critical-infrastructure surveillance network spans thousands of kilometers, with critical information scattered across unconnected sensors and surveillance systems. This fragmentation, along with human operation for threat detection and classification, makes the system suffer from high false-alarm rates, which constrain timely and reliable threat assessment.
Recent advancements in technology, AI, and Machine Learning offer new methods to enhance the system, such as multi-sensor fusion and AI analysis for automating surveillance systems, leaving only the critical decision-making part with manual operators. This survey provides a structured review of the five major domains of such an AI-based system that can potentially automate surveillance and reduce operators' fatigue and high false-alarm rates, providing accurate threat detection and classification. Our principal observation is that despite the growth in AI and ML, there exists no such framework that unifies diverse sensors and performs accurate threat detection and five-level threat classification across multiple domains. These gaps motivate the design and study of a unified, multi-agent, AI-driven surveillance platform tailored to Indian Defence requirements.
\end{abstract}

\vspace{\baselineskip}

\begin{IEEEkeywords}
Defence surveillance, multi-modal sensor fusion, multi-agent systems, threat detection and classification, explainable artificial intelligence, human-in-the-loop, adversarial robustness, edge computing, deep learning, literature survey
\end{IEEEkeywords}

\section{Introduction}
\label{sec:introduction}

India maintains active surveillance obligations along more than 15{,}000~km of land borders with Pakistan, Bangladesh, Myanmar, and Nepal, a 7{,}500~km coastline, and an expanding portfolio of critical infrastructure installations~\cite{indiaai2024springer}. Each segment employs its own constellation of radars, thermal imagers, unattended ground sensors, electro-optical cameras, and communication relays, generating data volumes that far exceed the analytical capacity of human operators~\cite{meng2023multisource}. Although individual sensor technologies have matured significantly over the years, the lack of a unified system that can connect and reason across multiple sensor types and geographic regions leaves today's surveillance infrastructure fundamentally fragile. A scientometric analysis of defence-technology publications between 2015 and 2025 confirms a fourfold increase in AI-related land-combat research, yet integrated multi-domain surveillance architectures remain conspicuously under-represented~\cite{emerging2025scientometric}.

Two critical problems make this gap even more dangerous. First, the high false alarm rates seen across deployed perimeter security systems place a growing burden on human operators, every false trigger pulls attention and resources away from threats that actually matter~\cite{deepsurvey2023anomaly}. Second, most existing alarm systems treat all alerts the same way, meaning a deer crossing a perimeter sensor and a coordinated military incursion receive the same level of urgency, making it impossible for operators to prioritise their response effectively~\cite{cctv2024anomaly}. Decision-support systems that incorporate explainability techniques such as SHAP values or attention-based saliency maps have shown encouraging results in civilian surveillance for identifying and triaging anomalies~\cite{zhang2022xai}, but translating these approaches into military environments, where decisions must be made faster, accountability is higher, and the consequences of errors are far greater, remains an unsolved challenge~\cite{military2024xai}.

Multi-agent AI architectures offer a principled response to both failures. Rather than replacing human judgment, a network of specialised software agents can function as a filtering and escalation layer that operationalises algorithmic situational awareness: each agent ingests a defined sensor stream, applies domain-specific inference, and feeds a threat-correlation agent that synthesises inputs into a graduated alert~\cite{hassan2025multiagent, agenticcyber2025}. Such architectures have demonstrated measurable gains in cyber-threat correlation and battlefield command support~\cite{zhao2022deepai}, and purpose-built human-autonomy teaming frameworks preserve commander authority over critical engagement decisions~\cite{autonomyteaming2024, meaningful2024hitl}. Extending this paradigm to physical-domain defence surveillance requires integration across five technical pillars that, to date, have been investigated largely in isolation.

This paper makes two contributions. First, it provides a structured literature review spanning multi-modal sensor fusion~\cite{meng2023multisource}, multi-agent threat detection~\cite{hassan2025multiagent, agenticcyber2025}, false-alarm reduction and explainable AI~\cite{zhang2022xai, deepsurvey2023anomaly}, geospatial awareness with human-in-the-loop control~\cite{karimipour2021iot, hitl2023trust}, and real-time distributed architectures with adversarial robustness~\cite{flink2024kafka, advyolo2022}. Second, it identifies and formalises the integration gap across these five domains as the primary bottleneck obstructing India's transition from sensor-rich but intelligence-poor surveillance to a unified, graduated threat-management grid. Section~II traces the historical and policy background. Sections~III through~VII examine each domain. Section~VIII presents a comparative taxonomy, analyses open gaps, outlines a proposed architectural direction, and concludes the paper.



\section{Background}
\label{sec:background}

India's post-independence approach towards managing its borders can be characterized by manpower-centric operations relying heavily on traditional techniques such as fences, patrolling, optical lookout posts, and high-frequency radio communications~\cite{indiaai2024springer}. Over the last several decades, the primary method for collecting information about the conditions at the frontline was through the eyes and ears of the soldiers stationed in border posts combined with basic surveillance tools such as binoculars, alarm wires, and searchlight systems installed on vehicles~\cite{indiaai2024springer}. Though a known approach, this technique does not scale effectively to India's size, variety of geographical regions, and harsh weather conditions and provides very little ability to conduct ongoing surveillance around-the-clock.

The Comprehensive Integrated Border Management System (CIBMS) may be considered as the largest breakthrough in terms of both technology and doctrine of Indian border surveillance. Designed to replace traditional physical barriers along the India-Pakistan and India-Bangladesh borders, CIBMS consists of various sensors ranging from thermal imaging cameras and radar systems mounted on tethered aerostats, to laser and infrared sensors that detect intrusion attempts and unattended ground sensors. Further refinements added various radars for detecting battlefield threats, human presence, and seismic detectors able to pick up sounds created during tunnel construction~\cite{indiaai2024springer}. Despite the transition to C2-centric surveillance, CIBMS' alert generation still relies on rule-based threshold logic and does not discriminate between false positives and actual violations of border integrity~\cite{cctv2024anomaly}.

There is a parallel revolution underway in the military surveillance R\&D sphere, with rule-based analysis giving way to machine learning and deep neural networks allowing for processing vast amounts of sensor data faster than even human analysts are able to~\cite{threatdetect2024rt}. Several convolutional architectures were shown to have good performance at detecting military vehicles and weapon systems~\cite{weapon2023yolov8} as well as military equipment visible in satellite imagery~\cite{satellite2024yolov8}. Combined CNN-LSTM architectures for processing sequential data improved anomaly detection accuracy and decreased false-positive rates~\cite{deepsurvey2023anomaly}. An example of successful implementation of AI pipelines is Project Maven launched by the U.S.\ Department of Defense, it showed that with machine learning assistance, 20 people were able to process the same amount of full-motion video footage as 2{,}000 people could without ML~\cite{cset2024maven, rand2023maven}.

The response from India's defence leadership has been quite proactive. The Ministry of Defence prioritized development of artificial intelligence solutions and increased CAIR's remit to include autonomous surveillance systems, target recognition, and AI-powered command systems~\cite{indiaai2024springer}. Commercial off-the-shelf AI providers and data fusion systems have previously been assessed for cross-database correlations and anomaly detection capability in intelligence analysis~\cite{cset2025ai}. Nevertheless, an analysis of India's defence-AI ecosystem reveals a prevalence of pilots, which are stovepipe within each branch of the defence forces due to insufficient interoperability capabilities and conceptually divergent data semantics~\cite{indiaai2024springer}.

There are several key concerns that drive this survey effort. One of the biggest problems is bandwidth limitations in contested border regions, which requires development of edge-computing solutions able to provide preliminary threat assessment at sensor nodes~\cite{flink2024kafka}. In addition, environmental noise sources, including sandstorms occurring near the border with Rajasthan and river level changes caused by seasonal monsoons in the border region with Bangladesh, generate long-term false alarm signals, which overwhelm threshold-based detection methods~\cite{deepsurvey2023anomaly}. Furthermore, deep learning classifiers suffer from poor interpretability, raising concerns about military accountability when making decisions related to escalation of force in accordance with international humanitarian laws~\cite{military2024xai, meaningful2024hitl}. Legacy sensor devices used in the CIBMS sector lack unified data interfaces, making real-time fusion with AI-enabled sensors difficult~\cite{meng2023multisource}. Finally, modern adversaries are able to attack object detection models via adversarial patches and GPS signal spoofing~\cite{advyolo2022, milai2025risk}. All of the mentioned concerns represent open challenges that this survey seeks to address, starting with a systematic overview of the five key pillars.

% NOTE: table* used for full-width layout; 5 content columns
%       cannot fit legibly within a single IEEEtran column.
\begin{table*}[!t]
\centering
\caption{Taxonomy of Defence Surveillance Literature}
\label{tab:taxonomy}
\renewcommand{\arraystretch}{1.3}
\footnotesize
\begin{tabular}{>{\raggedright\arraybackslash}p{1.5cm}>{\raggedright\arraybackslash}p{3.6cm}>{\raggedright\arraybackslash}p{2.4cm}>{\raggedright\arraybackslash}p{2.6cm}>{\raggedright\arraybackslash}p{3.8cm}}
\hline
\textbf{Domain} & \textbf{Representative Approaches} & \textbf{Sensor / Data Modality} & \textbf{Key Technique} & \textbf{Limitation Noted} \\
\hline
A: Sensor Fusion
  & Multi-source information fusion for C3I systems; attention-based cross-modal alignment
  & Radar, thermal, acoustic, optical, RF
  & Feature-level and decision-level multi-source fusion
  & Evaluated on pairwise modalities only; no adversarial degradation testing \\
\hline
B: Multi-Agent AI
  & Modular MAS with role-specialised sub-agents; YOLOv8 military object detection
  & Video, radar returns, satellite imagery, system logs
  & Agent decomposition with cross-source threat correlation
  & Homogeneous task distribution; no cross-modal agent orchestration \\
\hline
C: False Alarm / XAI
  & Deep-learning anomaly detection; SHAP/LIME explanation methods
  & Surveillance video, CCTV feeds, IDS telemetry
  & Temporal consistency filtering; post-hoc feature attribution
  & XAI latency not evaluated within real-time alert pipelines \\
\hline
D: Geospatial / HITL
  & GIS-based threat dashboards; trust calibration for autonomy teaming
  & GIS terrain layers, sensor metadata, force positions
  & Terrain-aware threat mapping; human authorisation loops
  & Tested in single-operator, single-threat scenarios only \\
\hline
E: Distributed / Adversarial
  & Kafka + Flink real-time streaming; adversarial patch attacks on YOLO detectors
  & IoT telemetry, edge video, RF uplinks
  & Edge preprocessing with stream-processing inference; adversarial training
  & No evaluation under simultaneous node failure and adversarial injection \\
\hline
\end{tabular}
\end{table*}
